\documentclass[11pt]{article}
\usepackage{amsmath,amsfonts}
\usepackage{graphicx}
%\usepackage{xcolor}
%\definecolor{gray}{rgb}{.75,.75,.75}
%\usepackage[landscape]{geometry}
%\usepackage{hyperref}
%\usepackage[bookmarks=true, pdfborder={0 0 .5}, urlbordercolor=gray]{hyperref}

\renewcommand{\thefootnote}{\fnsymbol{footnote}}


\addtolength{\topmargin}{-.4in}
\addtolength{\textheight}{.6in}
\addtolength{\oddsidemargin}{-.4in}
\addtolength{\textwidth}{.8in}
\pagestyle{empty}

\newcommand{\ds}{\displaystyle}
\newcommand{\ts}{\textstyle}

\newcommand{\Z}{\mathbb{Z}}
\newcommand{\R}{\mathbb{R}}
\newcommand{\C}{\mathbb{C}}
\newcommand{\EE}{\mathcal{E}}
\newcommand{\tZ}{\tilde{\Z}}

\newcommand{\sech}{\mathrm{sech}}
\newcommand{\csch}{\mathrm{csch}}


\newcommand{\Sym}{\mathrm{Sym}}

\renewcommand{\circle}[1]{{\bigcirc\hspace{-.75em}#1}}

\begin{document}

\footnote[0]{Notes by Peter Magyar \ {\tt magyar@math.msu.edu}}
\vspace{-1em}

\centerline{\bf Math 133\hspace{1.5em} \hfill{\Large\bf
Power Series}\hfill Stewart \S11.8}
\vspace{1em}

\noindent 
{\bf Series of functions.}
The main purpose of series is to write an interesting, 
complicated quantity as an infinite sum of simple quantities,
so that finite partial sums approximate the original quantity. 
For example, it is a fact (explained in \S11.9) that:
$$
\tfrac14 \pi \ =\ 1 -\tfrac1{3} + \tfrac15 -\tfrac17 + \tfrac19-\cdots\,,
$$
and we can approximate $\pi$ by taking enough terms of this series (and multiplying by 4).\footnote{
An irrational number like $\pi$ can never be written exactly as a fraction
or a finite decimal, so simple ways to approximate are the best we can hope for.
}
So far, we have no tools to find or prove such formulas.
Our Tests determine whether a series converges, 
but they tell us nothing about what it converges to,
nor how to write a given interesting quantity as a series.

We have looked only at series of {\it numbers} manipulated with basic algebra and limits. 
%The one case in which this
%gave a satisfactory answer was the geometric series:
%$$
%\sum_{n=0}^\infty c r^n \ =\ 
%\lim_{N\to\infty} \sum_{n=0}^N c r^n \ =\ 
%\lim_{N\to\infty} c\,\frac{\ 1{-}r^{N+1}}{1{-}r\ } \ =\
%\left\{\!\!\begin{array}{cl}
%c\,\dfrac{1}{1{-}r} & \text{ for $|r|<1$}\\[1em]
%\text{diverges} & \text{ for $|r|\geq 1$.}
%\end{array}\right.
%$$
In the next few sections, we will learn about series of {\it functions},
and use calculus to write interesting, complicated functions as 
infinite sums of simple functions.

\begin{quote}
{\it Definition:} A {\it power series} is a function of $x$ whose output
is the sum of an infinite series:
$$
f(x) \ =\ \sum_{n=0}^\infty c_n(x{-}a)^n
\ =\ c_0 + c_1 (x{-}a) + c_2 (x{-}a)^2 + c_3 (x{-}a)^3 + \cdots
$$
Here the numbers $c_0,c_1,c_2,\ldots$ are called the {\it coefficients},
and $a$ is called the {\it center}.
The domain of the function (the set of acceptable inputs) contains those
values of $x$ for which the series converges.
\end{quote}

\noindent
The simplest power series is:
$$
f(x) \ =\ \sum_{n=0}^\infty x^n \ =\ 1 + x+x^2 + x^3+\cdots .
$$
This is a function which takes any particular input $x=r$ 
to the output
$f(r)=\sum_{n=0}^\infty r^n = 1 + r + r^2 + r^3+\cdots$, 
the sum of the corresponding geometric series.
As we know, this converges exactly when $|r|<1$,
so the domain of $f(x)$ is $|x|<1$,
the interval $x\in (-1,1)$.
We also know a simple formula for this sum: 
$f(x) = \frac1{1-x}$, which means we can write 
this simple rational function 
as a kind of infinite polynomial function:
$$
f(x) \ =\ \frac{1}{1-x} \ =\ 1 + x+x^2 + x^3+\cdots  
\text{\ \ \ for $x\in(-1,1)$.}
$$
A general power series might not have a simple formula for the sum,
but any standard function can be written as a power series, as we shall see
in \S11.10.
\newpage

Taking a finite partial sum like $1+x+x^2$ gives an approximation
to $f(x)=\frac{1}{1-x}$ for each $x\in(-1,1)$, so their graphs are close to each other above this interval:
\\[1em]
\centerline{
\includegraphics[width=3in]
{11-8-Converge.png} 
}
\vspace{0em}

\noindent
The series converges fastest, the approximation is closest,
when $|x|$ is very small; that is, close to the center $x=0$.
The graph of $1+x+x^2$ is just a parabola,
but it is already a good approximation to $y=f(x)=\frac1{1-x}$ 
for $x\in(-\frac12,\frac12)$, and taking 
more terms, such as $1+x+x^2+x^3+x^4$, improves the approximation for $x\in(-1,1)$.
The approximation is least accurate close to the vertical asymptote $x=1$, since 
polynomial functions are always continuous.
The approximations are useless outside the interval of convergence.

We can modify $f(x)$ to get power series for related functions.
For example, if we substitute $1{-}x$ into the series for $f(x)$, 
we get a power series with center $x=1$:
$$\begin{array}{rcl}
\frac1x &=& \frac1{1-(1{-}x)} 
\ =\ 1 + (1{-}x) + (1{-}x)^2 + (1{-}x)^3 + \cdots
\\[.5em]
&=& 1 - (x{-}1) + (x{-}1)^2 - (x{-}1)^3 + \cdots\,.
\end{array}$$
Since the previous series $\sum_{n=0}^\infty x^n$ converges when $|x|<1$, 
the new series converges when $|1{-}x|<1$, 
the interval $x\in(0,2)$ centered at $x=1$. 
The partial sum approximations to $\frac1x$
are most accurate when $|1{-}x|$ is very
small, close to the center.
\\[1em]
{\bf Power Series Convergence Theorem.}
Any power series $f(x) \ =\ \sum_{n=0}^n c_n(x{-}a)^n$ has
one of three types of convergence:
\begin{itemize}
\item The series converges for all $x$.

\item The series converges for $|x{-}a|<R$, the interval $x\in(a{}-R,a{+}R)$, and it diverges for $|x{-}a|>R$,
where $R>0$ is a value called the {\it radius of convergence}.\footnote{ 
The convergence at the endpoints $x=a{-}R,a{+}R$ must be determined separately.
}

\item The series converges only at the center $x=a$ and diverges otherwise.
\end{itemize}

\noindent
Given a power series, we apply one of our Convergence Tests, usually the Ratio Test, to show which values of $x$ make the
series converge.

\addtolength{\textheight}{1in}

\newpage

\addtolength{\topmargin}{-.5in}

\noindent
{\sc example:} Determine the domain of convergence of 
\ $\ds\sum_{n=0}^\infty \frac{x^n}{n!}$.
\\
Since the factorial on the bottom grows much faster than the 
exponential function on top, we guess that this always converges.
To prove this, the Ratio Test is appropriate because the terms of the
series have a growing number of factors, most of which 
cancel:
$$
\frac{a_{n+1}}{a_n}
\ =\ 
\frac{x^{n+1}}{(n{+}1)!}\left/\frac{x^n}{n!}\right.
\ =\
\frac{x^{n+1}}{x^n}\cdot
\frac{n!}{(n{+}1)!} 
$$
$$
\ =\ 
x \cdot
\frac{n\,(n{-}1)(n{-}2)\cdots (2)(1)}
{(n{+}1)\,n\,(n{-}1)(n{-}2)\cdots (2)(1)}
\ =\
\frac{x}{n{+}1}\,.
$$
\\[0em]
Taking any particular value of $x$, even a large value like $x=1000$,
gives: 
$$
L\ =\ \lim\limits_{n\to\infty}\left|\frac{a_{n+1}}{a_n}\right| 
\ =\  \lim\limits_{n\to\infty}\frac{|x|}{n{+}1} \ =\  0\,.
$$
Since $L <1$, the sequence converges for all $x$.
\\[.5em]
{\sc example:} Determine the domain of convergence of \ \
$\ds\sum_{n=0}^\infty (-1)^n\frac{(x{-}2)^{2n+1}}{n^2\ln(n)}$\,.
\\
This is a power series with coefficients $c_n=\frac{(-1)^n}{n^2\ln(n)}$
and center $a=2$. Ratio Test:
$$
L\ =\ 
\lim_{n\to\infty}
\left|\frac{a_{n+1}}{a_n}\right|
\ =\ 
\lim_{n\to\infty}\left|
\frac{(x{-}2)^{2n+3}}{(n{+}1)^2\ln(n{+}1)}\left/
\frac{(x{-}2)^{2n+1}}{n^2\ln(n)}\right.\right|
$$
$$
\ =\
|x{-}2|^2\cdot\lim_{n\to\infty}
\frac{n^2}{(n{+}1)^2}
\cdot
\lim_{n\to\infty}\frac{\ln(n)}{\ln(n{+}1)}
\ =\ |x{-}2|^2\cdot 1\cdot 1 \ =\ |x{-}2|^2\,,
$$
where the last two limits are determined by L'H\^opital's Rule.
Thus, the series converges when $L = |x{-}2|^2 <1$, i.e.~when $|x{-}2|<1$,
and diverges otherwise.
The radius of convergence is $R=1$, and the open interval of convergence
is $x\in(1,3)$. We do not worry about convergence at the endpoints.
\\[.5em]
{\sc example:} Determine the domain of convergence of 
$\ds\sum_{n=0}^\infty \frac{(2x{+}4)^{3n}}{3^n}$\,.
\\
Factoring $2x{+}4=2(x{+}2)$, we can rewrite this as a power series:
$$
\sum_{n=0}^\infty \frac{(2x{+}4)^{3n}}{3^n}
\ =\
\sum_{n=0}^\infty 
\frac{2^{3n}}{3^n}\, (x{+}2)^{3n}
\ =\
\sum_{n=0}^\infty 
\left(\frac{8}{3}\right)^{\!\!n}\! (x{+}2)^{3n}
\,.
$$
Not every power $(x{+}2)^n$ appears, only multiple-of-3 powers $
a_n=(\frac83)^n(x{+}2)^{3n}$,
but this still fits into the definition of power series since
we can think of the missing terms as having zero coefficients.
The Ratio Test gives $L =\lim_{n\to\infty}|\frac{a_{n+1}}{a_n}|
= \frac83 |x{+}2|^3$, so the series converges when $\frac83 |x{+}2|^3< 1$,
i.e.~when $|x{+}2|<\sqrt[3]{3/8}=\frac12\sqrt[3]3$. 
The radius of convergence is $R=\frac12\sqrt[3]3$, 
and the center is the point where $x{+}2=0$, 
namely $x=-2$; so the open interval of convergence is: 
$$
x\,\in\,(-2-\tfrac12\sqrt[3]3,\,-2+\tfrac12\sqrt[3]3) \,\approx\, 
(-2.72,-1.29).
$$
{\sc example:} Find a power series $\sum_{n=0}^\infty c_n x^n$ which converges only at $x=0$.
We must choose coefficients $c_n$ which grow fast enough to overwhelm the decrease of 
$x^n$, no matter how small we take $x\neq 0$. The factorial $c_n=n!$ does it:
applying the Ratio Test to $\sum_{n=0}^\infty n!\,x^n$ gives
$L=\lim_{n\to \infty} (n{+}1)x=\infty$ for {\it any} value of $x\neq 0$,
showing divergence. Of course, for $x=0$, all the higher terms vanish, and the series converges.





\end{document}

%%%%%%%%%%%  Picture  %%%%%%%%%%
\\[0em]
\centerline{
%\includegraphics[width=2in]
{1-8-Discont-iv.png}
\qquad \qquad 
%\includegraphics[width=2in]
{1-8-Discont-v.png} 
}
\vspace{0em}

\noindent
%%%%%%%%%%%%%%%%%%%%%%%%%

%%%%%%%%%%%  Table  %%%%%%%%%%
In fact, we get the following derivatives:
$$\begin{array}{|c||c|c|c|c|c|c|}
\hline &&&&&& \\[-.9em] 
f(x) & \sin(x) & \cos(x) & \tan(x)  & \sec(x) & \csc(x) & \cot(x)
\\[.2em] \hline &&&&&& \\[-.9em]
f'(x) &\cos(x) & -\sin(x) & \sec^2(x) & \tan(x)\sec(x) & -\cot(x)\csc(x) & -\csc^2(x)
\\[.1em] \hline
\end{array}$$
\\[1em]
%%%%%%%%%%%%%%%%%%%%%%%%%%



















